% Owner: classification_e. Root and boundary_d full-source referee PASS.
\begin{proposition}[E.2.22 (Proved): propagation from a known open subset of a disc]
Let $\pi:\mathcal X\to\Delta$ be a proper holomorphic
submersion with fibres diffeomorphic to $S^6$.  Suppose
there is a nonempty open subset $U\subset\Delta$ such
that every $X_s$, $s\in U$, is a member of the fixed
coprime multiplicity family $(3k,4l)$ of E.2.1.
Then every fibre over $\Delta$ is a member of that
same family.
\end{proposition}
\begin{proof}
Local miniversality in E.2.1 supplies actual
holomorphic comparisons with the prescribed family
locally over $U$, with locally constant integral
markings.  Thus the known open set has the relative
root and normal-crossing charts used below; only
pointwise isomorphisms are not being substituted
for these relative comparisons.

The proof first propagates the actual geometric data
outside finitely many successive discrete analytic
exceptional sets.  It then fills these sets in reverse
order by E.2.21.  Later exceptional sets are only
claimed to be discrete on the domain then under
consideration; they may accumulate at points removed
at an earlier stage.

\emph{The intrinsic line and the relative pencil.}
Put $A=12kl$ and $D=kl-3k-4l\ne0$.
The argument of E.2.5 constructs the unique line
bundle $\mathcal L$ on the whole total space with
\[
 \mathcal L^D=K_{\mathcal X/\Delta}^{A}.
\]
Indeed Ehresmann transport over the contractible disc
gives $H^1(\mathcal X,\Z)=H^2(\mathcal X,\Z)=0$,
so the exponential sequence identifies its Picard
group with a complex vector space.  This construction
needs no hypothesis on the fibres outside $U$.

Proper coherence makes $\mathcal V=\pi_*\mathcal L$
a torsion-free, hence locally free, sheaf on the
disc.  Its rank is two because it has that rank
on $U$.  At every parameter its fibre injects
into $H^0(X_s,\mathcal L_s)$ by the multiplication
exact sequence used in E.2.5.  Choose a frame
of $\mathcal V$ on the disc.  The common zero
locus of its two evaluated sections is closed
analytic in $\mathcal X$.  Its image is closed
analytic by properness and Remmert's theorem
\cite[Chapter~II, Theorem~8.8]{Demailly2012E}.
It misses $U$, so is a discrete analytic subset.
Remove it temporarily.

On the remaining connected parameter curve $S$
the sections define a proper holomorphic map
\[
 g:\mathcal X_S\longrightarrow\mathbb P^1\times S,
 \qquad \mathcal L=g^*\mathcal O(1).
 \tag{E.2.22.1}
\]
Each $g_s$ is nonconstant because its two sections
are independent.  Its fibres have pure dimension
two; the smooth total space and smooth target
then give flatness.  The fibres are connected:
the finite Stein factor is generically of degree
one, as it is so above $U$, and a finite
birational map to the normal target is an
isomorphism.  Its source is irreducible, so
there is no additional vertical component
in the Stein factor.  We retain
\[
 K_{\mathcal X_S/S}^{A}=g^*\mathcal O(D).
 \tag{E.2.22.2}
\]

\emph{Critical values and ordinary tori.}
The relative critical locus of $g$ is analytic
and proper over the target.  Its image has
finite intersection with each target fibre:
the nonconstant proper map $g_s$ to a curve
has finitely many critical values.  Thus the
reduced critical image is finite over $S$.
Its one-dimensional part has degree three,
as can be read over $U$; every such component
dominates $S$.  Remove the images of its
zero-dimensional components and its branch
and singular points.  These are further
proper analytic discrete subsets, absent
over the known open set.

A finite etale parameter cover now labels
the three critical sections.  Pointwise
classification after this cover descends;
images of later discrete exceptional sets
remain discrete because this cover is finite.
One may work on each connected component,
each of which meets the inverse image of $U$.
Send the labelled sections to $0,1,\infty$.
There are exactly these three critical values
on the current parameter domain.

The ordinary family over
$B^\circ\times S$, where
$B^\circ=\mathbb P^1\setminus\{0,1,\infty\}$,
is smooth and proper.  Its compact surface
fibres have the topology of the known two-tori
by Ehresmann's theorem and connectedness.
Adjunction gives integral $c_1(K_F)=0$,
since $F$ is a smooth fibre in a complex
six-sphere and its normal line is trivial.
These surfaces are complex tori.  To spell
out the classification step, $b_1(F)=4$
is even, so $F$ is K\"ahler.  The zero
canonical class excludes exceptional curves
and ruled surfaces.  On a K\"ahler surface,
a nonzero effective divisor cannot have
zero real Chern class; hence a nonzero
pluricanonical section has no zeros and
cannot give positive Kodaira dimension.
The compact-surface classification leaves
the Kodaira-dimension-zero list, in which
$b_1=4$ identifies a two-torus
\cite[Chapter~VI, Theorem~1.1 and Table~10]{BHPV2004E}.

\emph{The finite root charts.}
For any labelled critical section $\Gamma_j$,
every component of the Cartier divisor $g^*\Gamma_j$
dominates $S$.  A vertical divisor component would
be an entire fibre $X_s$, making $g_s$ constant.
Every dominating component meets $U$, where the
reduced fibre is irreducible.  Thus there is
exactly one such component, and its multiplicity
is the known one.  In particular
\[
 g^*\Gamma_1=(3k)\mathcal D_1,
 \qquad g^*\Gamma_2=(4l)\mathcal D_2,
 \qquad g^*\Gamma_0=\mathcal W,
 \tag{E.2.22.3}
\]
with reduced total hypersurfaces on the right.
For clarity, two distinct dominating components
could not both restrict to the known irreducible
reduced fibre.  Each is a Cartier prime in
the smooth total space, is torsion-free and
hence flat over the smooth parameter curve,
and has pure two-dimensional fibres.  Thus
its fibre over any parameter in $U$ contains
the whole known irreducible reduced surface.
At a general smooth point of
that fibre the known local equation has only
one reduced total branch.  Properness forces
every dominating component to meet that open
part of the known fibre, a contradiction.

For a finite multiplicity $m$, write locally
$q=a^m u$, with $a$ a reduced local equation
of $\mathcal D_j$ and $u$ a unit.  The normalization
of the cyclic root base change $r^m=q$ is
locally the disjoint union of the graphs
\[
 r=\zeta a u^{1/m},\qquad \zeta^m=1.
\]
It is therefore smooth and etale over the
actual total space, with free cyclic deck
action.  On $r\ne0$ its map to the root
disc times $S$ is already smooth.  Its
critical locus is supported over $r=0$,
is proper over $S$, and is empty over $U$.
Delete its discrete analytic image for
each of the two roots.  The resulting
proper smooth root families have connected
fibres by their Stein factors and are
families of two-tori by the same topology,
adjunction and surface argument above.

\emph{A finite analytic test for relative normal crossings.}
The total cusp $\mathcal W$ in (E.2.22.3)
is a reduced embedded hypersurface in
the smooth fourfold.  Take its finite
analytic normalization
\[
 \nu:\widetilde{\mathcal W}\longrightarrow
       \mathcal W\longrightarrow\mathcal X_S.
\]
Existence and finiteness are the analytic
normalization theorem
\cite[Chapter~II, Definition~7.9 and Theorem~7.12]{Demailly2012E}.
The normalization is proper over $S$.
First remove the proper images of its
singular locus, and then of the critical
locus of its map to $S$.  Each is analytic
and misses $U$.  On the resulting domain
the normalization is smooth over $S$.
Next remove the proper image of the
rank-drop locus of
\[
 T_{\widetilde{\mathcal W}/S}
       \longrightarrow\nu^*T_{\mathcal X_S/S}.
\]
This locus too is analytic and absent over
$U$.  The finite map $\nu$ is now an immersion.

For $r=2,3,4$, let $R_r$ be the locus
of ordered, pairwise distinct preimages
in the $r$-fold fibre product of
$\widetilde{\mathcal W}$ over $\mathcal X_S$.
Its partial diagonals are both open and
closed: a finite immersion is locally
an embedding, so its diagonal is open,
and separatedness makes it closed.
Consequently $R_r$ is closed analytic
in the proper fibre product and is
proper over $S$.

At a point of $R_r$ the $r$ local sheets
have relative conormal lines in the same
rank-three bundle $T^*_{\mathcal X_S/S}$.
On $R_2$ and $R_3$, delete the proper
images of the rank-drop loci of these
two or three lines.  Delete also the
image of $R_4$.  These are closed
analytic images and miss $U$, hence
are discrete.  There are only finitely
many such deletions.

At a remaining point the total cusp
has at most three local sheets, each
embedded smoothly and submersive over
$S$, with independent vertical normal
covectors.  Since the \emph{total}
$\mathcal W$ is a reduced embedded
hypersurface, its local ideal is exactly
the product of the distinct local
equations of those sheets.  No finite
subring or seminormality assumption
is being used.  Independence of the
vertical normals makes these equations,
the parameter on $S$, and the remaining
coordinates a holomorphic coordinate
system.  Thus
\[
 q=x,\qquad q=xy,\qquad\hbox{or}\qquad q=xyz
 \tag{E.2.22.4}
\]
after absorbing a unit.  These are
actual relative normal-crossing charts;
in particular the remaining cusp fibres
are reduced.  This proof did not assume
that normalization commutes with base
change before obtaining the charts.

\emph{The normalized toric pair and its action.}
The smooth proper normalization and
its conductor strata have the same
topology, incidence and normal degrees
as over $U$.  A further finite etale
cover can label these finitely many
strata.  The argument of E.2.18 applies
to every fibre: the six conductor
curves are a hexagon of $(-1)$-curves,
three alternate contractions give a
triangle in $\mathbb P^2$, and the
normalized pair is $(\mathrm{dP}_6,D)$.
Its relative logarithmic tangent bundle
is the evaluation pullback of a rank-two
bundle on the parameter curve.

In this evaluation description the
conductor descent maps are holomorphic
endomorphisms of that rank-two bundle.
They are constant along the compact
conductor curves.  They equal the identity
over $U$, where both toric logarithmic
fields descend.  The holomorphic identity
theorem on the connected parameter curve
makes them the identity everywhere.
This is the descent proof of E.2.19;
it requires a nonempty known open set,
not specifically a punctured disc.
Grauert evaluation, commutation on the
ordinary tori, completeness and the
implicit extension of the two primitive
toric periods then give the actual
marked cusp torus action locally on
the current parameter curve.  The
integral monodromy and triple-point
characters are constant by transport
and are the prescribed ones.

\emph{Continuation of the full period kernel.}
Fix an ordinary target point $t_*$.
The integral local system over
$B^\circ\times S$ has constant marked
peripheral representation, up to a
choice of marking: its matrices are
transported from those over $U$.
Parameter monodromy commutes with the
base representation.  It therefore
preserves the cusp invariant lattice
\[
 \ker(A_0-I)=\Z\widehat w\oplus\Z\widehat\delta.
\]
The wedge of the two corresponding
holomorphic period vectors at $t_*$
has a well-defined zero locus on $S$:
changing its integral basis multiplies
the wedge by a nonzero determinant.
It is nonzero over $U$, so that zero
locus is discrete analytic.  Remove
it as one final exceptional layer.

On a connected universal cover of
the remaining parameter curve, normalize
the Lie frame at $t_*$ by these two
columns.  The $\widehat\gamma$ column
then defines a holomorphic function
$c(s)$ by subtracting the fixed
reference value of $\beta^{\rm part}(t_*)$.
All other column identities at $t_*$
hold over a lift of $U$ and hence
everywhere by the identity theorem.

Now work on the universal cover of
the whole ordinary base times this
parameter cover.  The actual holomorphic
period kernel is a rank-two subbundle
of the complexified marked rank-four
lattice.  Compare it with
\[
 \ker\Pi_{c(s)}(t).
 \tag{E.2.22.5}
\]
The latter matrix always has complex
rank two, since its $\widehat w$ and
$\widehat\delta$ columns are the
standard basis.  These two holomorphic
Grassmannian maps agree over the
product with a lift of $U$.  They
therefore agree on the entire
connected ordinary cover.
This comparison does not assume
that the two actual columns were
independent at every $t$; equality
of the kernels proves it afterwards.
In a local actual Lie frame the
identity says that the full period
matrix is $W(t,s)\Pi_{c(s)}(t)$,
where $W$ is its two marked columns.
If $W$ had rank at most one, every
period column would have that rank,
contradicting the real rank four
of the compact torus lattice.

The actual ordinary and finite-root
fibres are compact tori.  The exact
period determinant is nonzero there.
Its sign is the negative sign on
$U$ and stays negative by connectedness
and nonvanishing.  Equality of the
period kernels also extends to the
smooth finite root charts; their
nonzero determinant preserves this
sign at the root point where the
maximum in R3 is attained.  Thus
$\operatorname{Im}c(s)+M<0$ by
the precise threshold of R3.
This places $c(s)$ in the
admissible interior domain, including
strictness at its boundary.  The
cusp multipliers are the corresponding
marked reference multipliers.  The
proof of E.2.20 now reconstructs
the actual cusp at each parameter:
relative normal-crossing transport
gives its known cover and incidence,
the action gives the same characters,
and adjacent-origin comparison gives
the monomial times unit extensions.
The proof works along any path from
$U$ and locally in the current
parameter; it does not need an
already known punctured neighbourhood
at the point being classified.

The finite affine lattice classes
and integral collar classes are
locally constant on this connected
domain, hence remain those over
$U$.  The additive comparison in
E.1 and E.2.21 identifies every
remaining actual threefold with
the reference member at $c(s)$.
The pointwise conclusion descends
from every finite parameter cover
used to label the critical or
conductor strata.

\emph{Filling the exceptional layers.}
We have made only finitely many
successive deletions.  At each stage
the deleted subset was closed
analytic and discrete on the
parameter domain then present.
Finite covers do not change that
property after taking their proper
images.  On the final domain all
fibres are known members of the
fixed pair.  At a point of the
last deleted set, choose a small
disc containing no other point of
that set and avoiding all earlier
deleted sets.  Its punctured disc
consists of known members.  E.2.21
identifies the central fibre as
well.  Fill the entire last layer
in this way, then the preceding
one, and continue in reverse order.
This is valid even when a later
layer accumulates at an earlier
deleted point.  Finally every
fibre over the original disc is
known.  No analytic curve selection
for an arbitrary open locus has
been assumed.
\end{proof}
